
\paragraph{\textbf{Lemma 2.}}
\SGARefTarget{sga3:xiv:7:lemma2}
{\itshape One has \(f'\equiv f\pmod W\). Conversely, every
\(f'\in E\) satisfying this condition arises from a maximal torus
\(T'\) of \(G\).}

Put \(\bar T=T_{\bar k}\) and \(\bar T'=T'_{\bar k}\), and choose
\(y\in G(\bar k)\) such that
\[
  y\bar T y^{-1}=\bar T'.
\]
Since \(y^q\bar T y^{-q}=\bar T'\), the element
\[
  n=y^{-1}y^q
\]
belongs to the normalizer \(N(T)(\bar k)\). The action of \(f'\)
on the points of \(\bar T\) is
\[
 \begin{aligned}
  f'(t)
    &=y^{-1}(yty^{-1})^q y\\
    &=y^{-1}y^q\,t^q\,y^{-q}y
      =nt^qn^{-1}.
 \end{aligned}
\]
If \(w\in W\) is represented by \(n\), then \(f'\) and
\(f\circ w\) have the same action on \(\bar T(\bar k)\), hence
on its characters. Thus \(f'=f\circ w\), and
\(f'\equiv f\pmod W\).

Conversely, represent a given \(w\in W\) by
\(n\in N(T)(\bar k)\). Lang's theorem gives
\[
  n=y^{-1}y^q
\]
for some \(y\in G(\bar k)\). Then
\(y\bar T y^{-1}\) is defined over \(k\), and the preceding
calculation shows that its corresponding element is
\(f\circ w\).

\paragraph{\textbf{Lemma 3.}}
\SGARefTarget{sga3:xiv:7:lemma3}
{\itshape Let \(X,R,W,E\) be as above, with \(R\) irreducible,
and let \(\varphi\in E/W\). There exist a representative
\(f\in E\) of \(\varphi\) and roots
\(\theta_1,\ldots,\theta_n\) such that:
\begin{enumerate}
\item \((\theta_1,\ldots,\theta_n)\) is a basis of \(X\);
\item \(R\) is the union of the orbits of the \(\theta_i\) under
the powers of \(f\): every \(a\in R\) has the form
\(a=f^m\theta_i\) for suitable \(m\) and \(i\).
\end{enumerate}}

The proof will be given below. We first need a lemma of linear
algebra.

\paragraph{\textbf{Lemma 4.}}
\SGARefTarget{sga3:xiv:7:lemma4}
{\itshape Let \(V\) be a vector space over \(k\), and let
\((\theta_1,\ldots,\theta_n)\) be a basis of the dual of
\(V\otimes_k\bar k\). Then there is an \(x\in V\) such that
\(\theta_i(x)\ne0\) for every \(i\).}

Let \(V^*\) be the dual of \(V\), and let \(W_i\) be the
\(\bar k\)-subspace of \(V^*\otimes_k\bar k\) generated by the
Galois conjugates of \(\theta_i\). It descends to a
\(k\)-subspace, again denoted \(W_i\). The natural map
\[
  W_1\otimes\cdots\otimes W_n
     \longrightarrow\bigwedge\nolimits^n V^*
\]
is nonzero: after extension to \(\bar k\), it contains
\(\theta_1\wedge\cdots\wedge\theta_n\ne0\). Hence \(V^*\) has a
basis \(u_1,\ldots,u_n\) with \(u_i\in W_i\). Choose
\(x\in V\) with \(u_i(x)\ne0\) for all \(i\), for example with
\(u_i(x)=1\). If \(\theta_i(x)=0\), every Galois conjugate of
\(\theta_i\) vanishes at \(x\), and so does \(u_i\), a
contradiction.

\paragraph{\textbf{End of the proof of the theorem.}}

By \SGARefLink{sga3:xiv:7:lemma2}{Lemmas 2} and 3, choose a maximal torus \(T\) whose element
\(f\) has the properties of \SGARefLink{sga3:xiv:7:lemma3}{Lemma 3}. Let \(Y\) be the dual of
\(X\). The Lie algebra of \(T\) satisfies
\[
  \mathfrak t(\bar k)\simeq Y\otimes_{\mathbb Z}\bar k,
\]
compatibly with the Galois action. Put
\(V=\mathfrak t(k)\). An element \(x\in V\) is regular precisely
when it is annihilated by no root \(\alpha\in R\), viewed as a
linear form in \(V^*\otimes_k\bar k\).

\SGARefLink{sga3:xiv:7:lemma4}{Lemma 4} supplies \(x\in V\) annihilated by none of the
\(\theta_i\). Every root is a Galois conjugate of one of the
\(\theta_i\), by property \textup{\SGARefLink{sga3:xiv:7:lemma2}{(2)}} of \SGARefLink{sga3:xiv:7:lemma3}{Lemma 3}. Hence no root
annihilates \(x\), and \(x\) is regular. This proves the theorem.
